\documentclass[aspectratio=169,11pt]{beamer}
% Shared Beamer style for practice contest solution walkthroughs.
\usepackage[utf8]{inputenc}
\usepackage[T1]{fontenc}
\usepackage{lmodern}
\usepackage{amsmath,amssymb}
\usepackage{xcolor}
\usepackage{hyperref}
\usepackage{listings}

\definecolor{CPNavy}{HTML}{172033}
\definecolor{CPBlue}{HTML}{2563EB}
\definecolor{CPGreen}{HTML}{16A34A}
\definecolor{CPOrange}{HTML}{F97316}
\definecolor{CPGray}{HTML}{64748B}
\definecolor{CPLight}{HTML}{F8FAFC}
\definecolor{CPCodeBg}{HTML}{F1F5F9}
\definecolor{CPCodeRule}{HTML}{CBD5E1}

\usetheme{default}
\usefonttheme{professionalfonts}
\setbeamertemplate{navigation symbols}{}
\setbeamersize{text margin left=0.65cm,text margin right=0.65cm}
\setbeamertemplate{itemize item}{\color{CPBlue}$\blacktriangleright$}
\setbeamertemplate{itemize subitem}{\color{CPGray}$\triangleright$}

\setbeamercolor{normal text}{fg=CPNavy,bg=white}
\setbeamercolor{frametitle}{fg=white,bg=CPNavy}
\setbeamercolor{title}{fg=white,bg=CPNavy}
\setbeamercolor{subtitle}{fg=CPLight,bg=CPNavy}
\hypersetup{colorlinks=true,urlcolor=CPBlue}

\setbeamertemplate{frametitle}{%
  \nointerlineskip%
  \begin{beamercolorbox}[wd=\paperwidth,ht=0.88cm,dp=0.22cm,leftskip=0.55cm,rightskip=0.35cm]{frametitle}
    \usebeamerfont{frametitle}\insertframetitle\hfill\small\insertframenumber/\inserttotalframenumber
  \end{beamercolorbox}%
}

\setbeamertemplate{footline}{%
  \leavevmode%
  \hbox{%
    \begin{beamercolorbox}[wd=.58\paperwidth,ht=2.4ex,dp=1ex,leftskip=.5cm]{author in head/foot}%
      \color{CPGray}\insertshorttitle
    \end{beamercolorbox}%
    \begin{beamercolorbox}[wd=.42\paperwidth,ht=2.4ex,dp=1ex,rightskip=.5cm plus1fil]{date in head/foot}%
      \hfill\color{CPGray}\insertshortauthor
    \end{beamercolorbox}%
  }%
}

\setbeamertemplate{title page}{%
  \vspace*{1.25cm}
  \begin{beamercolorbox}[wd=\paperwidth,sep=0.65cm,leftskip=0.15cm,rightskip=0.15cm]{title}
    {\color{white}\Huge\bfseries\inserttitle\par}
    \vspace{0.35cm}
    {\color{CPLight}\Large\insertsubtitle\par}
    \vspace{0.6cm}
    {\color{CPLight}\large\insertauthor\par}
  \end{beamercolorbox}
  \vfill
}

\newcommand{\sectiontag}[1]{\textbf{\color{CPBlue}#1}}
\newenvironment{tightitemize}{\begin{itemize}\small\setlength{\itemsep}{0.18cm}}{\end{itemize}}
\lstdefinestyle{cpcode}{
  language=Python,
  basicstyle=\ttfamily\tiny,
  keywordstyle=\color{CPBlue}\bfseries,
  commentstyle=\color{CPGray}\itshape,
  stringstyle=\color{CPGreen},
  backgroundcolor=\color{CPCodeBg},
  frame=single,
  framerule=0.25pt,
  rulecolor=\color{CPCodeRule},
  showstringspaces=false,
  columns=fullflexible,
  keepspaces=true,
  breaklines=true,
  breakatwhitespace=false,
  tabsize=4,
  xleftmargin=0.08cm,
  xrightmargin=0.08cm,
  aboveskip=0.08cm,
  belowskip=0.08cm
}


\title[dobra]{Dobra}
\subtitle{Day 2 solution walkthrough}
\author[Solution Deck]{Competitive Programming Bootcamp}
\date{}

\begin{document}

\begin{frame}[plain]
  \titlepage
\end{frame}

% BEGIN GENERATED STATEMENT INTERPRETATION
\begin{frame}{Interpret the Word Problem}
\small
\sectiontag{Statement excerpts:}
\begin{tightitemize}
  \item \textit{``do not contain \$3\$ sequential vowels''}
  \item \textit{``contain at least one letter''}
\end{tightitemize}

\vspace{0.18cm}
\sectiontag{Programming interpretation:}
\begin{tightitemize}
  \item Underscores are unknown letters that must be substituted.
  \item The recursive state needs only the position, consecutive vowel count, consecutive consonant count, and whether L appeared.
  \item Group choices into vowel, consonant-not-L, and L instead of trying all 26 letters separately.
\end{tightitemize}
\end{frame}
% END GENERATED STATEMENT INTERPRETATION







\begin{frame}{Problem Model}
\small
\sectiontag{Textbook connection:} CP4 3.2.2 Recursive Complete Search

\vspace{0.25cm}
\sectiontag{Core technique:} Grouped recursive backtracking

\vspace{0.25cm}
\begin{tightitemize}
  \item Underscores can become vowels, consonants other than L, or the letter L.
  \item A valid word has no three consecutive vowels or consonants.
  \item A valid word must contain at least one L.
\end{tightitemize}
\end{frame}

\begin{frame}{How to Recognize the Approach}
\begin{tightitemize}
  \item Trying all 26 letters per underscore is wasteful.
  \item The rules depend only on vowel/consonant category and whether L has appeared.
  \item Grouping choices gives the same count with far fewer branches.
\end{tightitemize}
\end{frame}

\begin{frame}{Algorithm}
\begin{tightitemize}
  \item Define backtrack(pos, vowels, consonants, has\_l).
  \item Prune immediately if vowels >= 3 or consonants >= 3.
  \item At the end, return 1 only if has\_l is true.
  \item For an underscore, branch into 5 vowels, 20 non-L consonants, and 1 L.
  \item For fixed letters, update the appropriate consecutive count.
\end{tightitemize}
\end{frame}

\begin{frame}{Walkthrough of the Key State}
\begin{tightitemize}
  \item The multiplier 5 counts all vowel replacements without listing them.
  \item The multiplier 20 counts consonants except L.
  \item The L branch sets has\_l to true and counts as a consonant.
\end{tightitemize}
\end{frame}

\begin{frame}{Why the Algorithm Is Correct}
\begin{tightitemize}
  \item Every concrete replacement belongs to exactly one of the three underscore groups.
  \item The state records all information needed for future validity checks.
  \item Invalid prefixes are pruned only after they already violate a problem rule, so no valid word is lost.
\end{tightitemize}
\end{frame}

\begin{frame}{Complexity and Implementation Notes}
\small
\sectiontag{Complexity:} O(3\textasciicircum{}U) recursive branches after grouping, where U is the number of underscores; O(L) recursion depth.

\vspace{0.3cm}
\begin{tightitemize}
  \item Memoization is possible but not required for this version.
  \item The branch weights are counts of letters, not probabilities.
\end{tightitemize}
\end{frame}



% BEGIN GENERATED CODE WALKTHROUGH
\begin{frame}[fragile]{Code: Input and State Function}
\scriptsize
\sectiontag{Source lines:} \texttt{dobra.py:38--45}

\vspace{0.08cm}
\begin{columns}[T,totalwidth=\textwidth]
\begin{column}{0.58\textwidth}
\begin{lstlisting}[style=cpcode]
import sys

word = sys.stdin.readline().strip()

VOWELS = "AEIOU"

def backtrack(pos, vowels, consonants, has_l):
\end{lstlisting}
\end{column}
\begin{column}{0.38\textwidth}
\sectiontag{What this chunk does}
\begin{tightitemize}
  \item The whole word is read once and shared by the recursive function.
  \item VOWELS defines the only letters that reset the consonant streak.
  \item The parameters describe everything needed to continue the search.
\end{tightitemize}
\end{column}
\end{columns}
\end{frame}

\begin{frame}[fragile]{Code: Prune and Finish}
\scriptsize
\sectiontag{Source lines:} \texttt{dobra.py:46--54}

\vspace{0.08cm}
\begin{columns}[T,totalwidth=\textwidth]
\begin{column}{0.58\textwidth}
\begin{lstlisting}[style=cpcode]
if vowels >= 3 or consonants >= 3:
    return 0

if pos == len(word):
    return 1 if has_l else 0

char = word[pos]
\end{lstlisting}
\end{column}
\begin{column}{0.38\textwidth}
\sectiontag{What this chunk does}
\begin{tightitemize}
  \item Any branch with three vowels or consonants in a row is immediately invalid.
  \item At the end of the word, the branch counts only if an L has appeared.
  \item The current character determines which transition to apply next.
\end{tightitemize}
\end{column}
\end{columns}
\end{frame}

\begin{frame}[fragile]{Code: Expand an Underscore}
\scriptsize
\sectiontag{Source lines:} \texttt{dobra.py:56--68}

\vspace{0.08cm}
\begin{columns}[T,totalwidth=\textwidth]
\begin{column}{0.58\textwidth}
\begin{lstlisting}[style=cpcode]
if char == "_":
    total = 0

    total += 5 * backtrack(pos + 1, vowels + 1, 0, has_l)

    total += 20 * backtrack(pos + 1, 0, consonants + 1, has_l)

    total += backtrack(pos + 1, 0, consonants + 1, True)

    return total
\end{lstlisting}
\end{column}
\begin{column}{0.38\textwidth}
\sectiontag{What this chunk does}
\begin{tightitemize}
  \item Five vowel choices have identical future behavior, so multiply that branch by 5.
  \item Twenty consonants other than L also share one transition.
  \item The single L branch turns on has\_l and counts as a consonant.
\end{tightitemize}
\end{column}
\end{columns}
\end{frame}

\begin{frame}[fragile]{Code: Handle Fixed Letters}
\scriptsize
\sectiontag{Source lines:} \texttt{dobra.py:70--79}

\vspace{0.08cm}
\begin{columns}[T,totalwidth=\textwidth]
\begin{column}{0.58\textwidth}
\begin{lstlisting}[style=cpcode]
    elif char in VOWELS:
        return backtrack(pos + 1, vowels + 1, 0, has_l)

    else:
        return backtrack(pos + 1, 0, consonants + 1, has_l or char == "L")

print(backtrack(0, 0, 0, False))
\end{lstlisting}
\end{column}
\begin{column}{0.38\textwidth}
\sectiontag{What this chunk does}
\begin{tightitemize}
  \item A fixed vowel increments the vowel streak and resets consonants.
  \item A fixed consonant increments the consonant streak and resets vowels.
  \item The final print starts the recursion with no letters processed and no L seen.
\end{tightitemize}
\end{column}
\end{columns}
\end{frame}
% END GENERATED CODE WALKTHROUGH

\begin{frame}{Source}
\small
\sectiontag{Solution file:} \texttt{practice\_contests/day\_2/dobra.py}

\vspace{0.3cm}
\sectiontag{Problem:} \url{https://open.kattis.com/problems/dobra}

\vspace{0.45cm}
Use this deck with the source code open beside it. The slides explain the reasoning; the code shows the exact input handling and output formatting.
\end{frame}

\end{document}
